%% kapa: exam=1 round=35 year=2024 period=2 form=A subject=appraisal_law subject_ko=감정평가관계법규 seq=33 number=33
%% source: Q-Net 감정평가사 기출문제 (cst00309) artlSeq=5213229 file=raw/1008/5213229/2교시 문제지 원본.pdf
%% conversion: kordoc 4.24.51 -> tools/gen1.py (auto); review=proofread-sonnet; flags=none
\begin{question}{33}
부동산등기법령상 권리에 관한 등기에 관한 설명으로 옳은 것은?
\begin{choices}
  \item 임차권을 정지조건부로 설정하는 청구권을 보전하려는 경우에도 가등기를 할 수 있다.
  \item 등기관이 등기를 마친 후 그 등기가 신청할 권한이 없는 자가 신청한 것임을 발견한 때에는 등기를 직권말소한다는 뜻을 통지하여야 한다.
  \item 환매특약등기는 이해관계 있는 제3자의 승낙이 없는 경우 부기로 할 수 없다.
  \item 미등기의 토지에 대해 매매계약서에 의하여 소유권을 증명하는 자는 그 토지에 관한 소유권보존등기를 신청할 수 있다.
  \item 등기관이 직권으로 등기를 말소한 처분에 대하여 관할 법원에 이의를 신청하면 등기 말소처분은 효력이 정지된다.
\end{choices}
\end{question}
\answer{1}
